\chapter{Glaucoma Symptoms}

\section*{Definition}
Glaucoma comprises a group of optic neuropathies characterized by progressive retinal ganglion cell death and characteristic optic nerve cupping, often associated with elevated intraocular pressure (IOP). The two major forms are primary open-angle glaucoma (POAG, insidious) and primary angle-closure glaucoma (PACG, acute).

\section*{What Happens in Your Body}
Elevated IOP (or pressure-independent mechanisms) causes mechanical stress and vascular compromise at the lamina cribrosa, disrupting axonal transport and leading to apoptotic retinal ganglion cell death. This results in progressive optic nerve head cupping and corresponding visual field loss that begins in the periphery.

\section*{Causes}
\begin{itemize}
  \item \textbf{Primary open-angle:} Idiopathic with risk factors including age > 40, African ancestry, family history, and myopia
  \item \textbf{Primary angle-closure:} Anatomically narrow angle with acute pupillary block
  \item \textbf{Pseudoexfoliation:} Fibrillar protein deposits obstructing trabecular meshwork
  \item \textbf{Pigmentary:} Iris pigment dispersion causing trabecular meshwork obstruction
  \item \textbf{Secondary:} Uveitic, neovascular (diabetic/ischemic), traumatic, or steroid-induced
\end{itemize}

\section*{When to See a Doctor}
See your doctor if:
\begin{itemize}
  \item Acute onset of severe unilateral eye pain, headache, nausea, and blurred vision (PACG)
  \item Colored halos around lights with fixed, mid-dilated pupil and cloudy cornea
  \item Incidental finding of elevated IOP or enlarged cup-to-disc ratio on routine exam
  \item Progressive loss of peripheral vision noted by you
  \item History of glaucoma in a first-degree relative (screening recommended)
\end{itemize}

\section*{Related Symptoms}
\begin{itemize}
  \item Gradual peripheral vision loss (tunnel vision in advanced POAG)
  \item Ocular pain and headache (more prominent in PACG)
  \item Nausea and vomiting (severe in acute angle-closure)
  \item Blurred vision and halos around lights
  \item Conjunctival injection and corneal edema in acute presentation
\end{itemize}
\textbf{Important:} Chronic open-angle glaucoma is typically asymptomatic until advanced, making regular screening essential; by the time vision loss is noticed, up to 40\% of retinal ganglion cells may already be lost.

\section*{Self-Care / Home Management}
\begin{itemize}
  \item Strict adherence to prescribed topical medications
  \item Technique optimization for eyedrop instillation (punctal occlusion to reduce systemic absorption)
  \item Avoidance of activities that spike IOP (weightlifting with breath-holding, inverted yoga poses)
  \item Sleeping with head elevated 30\textdegree{} may slightly lower nocturnal IOP
  \item Maintaining stable hydration (avoid drinking > 1 liter within 20 minutes)
\end{itemize}

\section*{How Your Doctor Will Evaluate You}
\begin{itemize}
  \item Goldmann applanation tonometry for IOP measurement
  \item Gonioscopy to assess anterior chamber angle configuration
  \item Slit-lamp examination of optic nerve head (cup-to-disc ratio, rim thinning, disc hemorrhage)
  \item Optical coherence tomography (OCT) of retinal nerve fiber layer and ganglion cell complex
  \item Standard automated perimetry (Humphrey 24-2 SITA) for visual field defect detection
  \item Central corneal thickness measurement (pachymetry) to adjust IOP interpretation
\end{itemize}

\section*{Treatment Options}
\begin{itemize}
  \item First-line medical: Prostaglandin analogs (latanoprost, travoprost) once daily
  \item Add-on therapy: Beta-blockers (timolol), alpha-2 agonists (brimonidine), carbonic anhydrase inhibitors (dorzolamide), or rho kinase inhibitors (netarsudil)
  \item Laser treatment: Selective laser trabeculoplasty (SLT) for POAG; laser peripheral iridotomy (LPI) for PACG
  \item Incisional surgery: Trabeculectomy with mitomycin C or glaucoma drainage implant
  \item Minimally invasive glaucoma surgery (MIGS: iStent, Hydrus, Kahook Dual Blade) combined with cataract extraction
  \item Cyclophotocoagulation (diode or micropulse) for refractory cases
\end{itemize}

\clearpage
